% Einheit 1 von 4 – Terme aufstellen und berechnen (bank/terme/e1.jsonl, zusammenbau v0.9)
\einheitenkopf[e1]{Einheit 1 von 4 · Terme aufstellen und berechnen}
\zweigzeile{Hier lernst du, Terme aufzustellen und zu berechnen – Variable als Platzhalter, Termwert, Term aus Sachtext und Figur · ab Kl. 6, je nach Buch bis Kl. 7 · P10 ×5 · baut auf: Addieren und Subtrahieren negativer Zahlen, Multiplizieren mit Vorzeichen, Dezimalzahlen und einfache Brüche als Vorzahlen}
\setcounter{aufgabe}{5}


% leiter: terme-e1-k1-s1-v1, terme-e1-k1-s1-v2, terme-e1-k1-s1-v3, terme-e1-k1-s2-v1
\begin{kbaufgabe}{Ich kann einen Term aufstellen.}
\begin{kbblock}
\kbanweisung{Kreuze den Term an, der dazu passt.}
\lbpaar{\kbteil{a)}{Das Vierfache einer Zahl $x$ wird um 3 vermindert.}{}
\kbkreuz{$4x - 3$}
\kbkreuz{$4 \cdot (x - 3)$}
\kbkreuz{$3 - 4x$}
}{\kbteil{b)}{Das Vierfache einer Zahl $x$ wird um 5 vermindert.}{}
\kbkreuz{$5 - 4x$}
\kbkreuz{$4x - 5$}
\kbkreuz{$4 \cdot (x - 5)$}}
\end{kbblock}
\begin{kbblock}
\lbpaar{\kbteil{c)}{Das Vierfache einer Zahl $x$ wird um 6 vermindert.}{}
\kbkreuz{$4 \cdot (x - 6)$}
\kbkreuz{$6 - 4x$}
\kbkreuz{$4x - 6$}
}{\item[]}
\end{kbblock}
\begin{kbblock}
\kbteil{d)}{Das Dreifache einer Zahl $x$ wird um 8 vergrößert. Schreibe dazu einen Term.}{}
\kbantwort{\kbfeld}
\end{kbblock}
\end{kbaufgabe}


% leiter: terme-e1-k1-s3-v1, terme-e1-k1-s4-v1, terme-e1-k1-s5-v1, terme-e1-k1-s6-v1
\begin{kbaufgabe}{Ich kann einen Term aufstellen. – weiter}
\begin{kbblock}
\kbteil{a)}{Die Figur zeigt ein Rechteck. Die Maße sind in cm angegeben. Schreibe einen Term für den Flächeninhalt des Rechtecks.}{}
\kbdarunter{\viereck[seiten={5,x,5,x}]{(0,0)}{(5,0)}{(5,2.5)}{(0,2.5)}}
\kbantwort{\kbfeld}
\end{kbblock}
\begin{kbblock}
\kbteil{b)}{Ein Heft kostet $a$ Euro, ein Stift kostet $b$ Euro. Lena kauft 4 Hefte und 3 Stifte. Schreibe einen Term für den Preis in Euro, den Lena zahlt.}{}
\kbantwort{\kbfeld}
\end{kbblock}
\begin{kbblock}
\kbteil{c)}{Ein Rechteck hat die Seiten $x$ cm und 9 cm. Schreibe einen Term für den Umfang. Fasse den Term zusammen.}{}
\kbantwort{\kbfeld}
\end{kbblock}
\begin{kbblock}
\kbteil{d)}{Eine Rechnung geht so: Verdreifache eine Zahl $x$, addiere 2 und verdopple die Summe. Schreibe dazu einen Term.}{}
\kbantwort{\kbfeld}
\end{kbblock}
\end{kbaufgabe}


% pruefung: terme-e1-k1-s7-v1, terme-e1-k1-s7-v3
\begin{kbaufgabe}{Ich finde auch in Aufgaben aus der Prüfung den passenden Term.}
\begin{kbblock}
\kbanweisung{Kreuze den passenden Term an.}
\kbteil{a)}{Die Differenz aus dem Dreifachen einer Zahl $x$ und 7 wird verdoppelt.}{P10 ’23}
\kbkreuz{$(3x - 7) : 2$}
\kbkreuz{$3x : 7 \cdot 2$}
\kbkreuz{$2 \cdot (3x - 7)$}
\kbkreuz{$3x - 7 \cdot 2$}
\end{kbblock}
\begin{kbblock}
\kbteil{b)}{Das Sechsfache einer Zahl $x$ wird um zwei vermindert.}{P10 ’17}
\kbkreuz{$6 \cdot (x - 2)$}
\kbkreuz{$2 - 6x$}
\kbkreuz{$6x - 2$}
\kbkreuz{$6x - 2x$}
\end{kbblock}
\end{kbaufgabe}


% ohne: terme-e1-k2-s1-v1, terme-e1-k2-s1-v2, terme-e1-k2-s1-v3
\begin{kbaufgabe}{Ich kann Termwerte berechnen.}
\begin{kbblock}
\kbanweisung{Berechne den Wert des Terms.}
\lbpaar{\kbteil{a)}{$4x - 3$ für $x = 5$ \quad \lbfeld}{}
}{\kbteil{b)}{$2x + 9$ für $x = -3$ \quad \lbfeld}{}}
\end{kbblock}
\begin{kbblock}
\lbpaar{\kbteil{c)}{$10 - 3x$ für $x = -2$ \quad \lbfeld}{}
}{\item[]}
\end{kbblock}
\end{kbaufgabe}


% ohne: terme-e1-k3-s1-v1, terme-e1-k3-s1-v2, terme-e1-k3-s1-v3
\begin{kbaufgabe}{Ich kann zu einem Term eine passende Situation erfinden.}
\begin{kbblock}
\kbanweisung{Erfinde eine Situation aus dem Alltag, zu der der Term … passt. Schreibe auf, wofür $x$ steht. Berechne den Term für … und erkläre, was das Ergebnis bedeutet.}
\kbteil{a)}{$15 + 3x$, $x = 2$}{}
\item[]\kbraum{2}
\end{kbblock}
\begin{kbblock}
\kbteil{b)}{$20 - 2x$, $x = 4$}{}
\item[]\kbraum{2}
\end{kbblock}
\begin{kbblock}
\kbteil{c)}{$10 + 4x$, $x = 3$}{}
\item[]\kbraum{2}
\end{kbblock}
\end{kbaufgabe}


% pflicht: terme-e1-k4-s1-v1, terme-e1-k4-s1-v2, terme-e1-k4-s1-v3
\begin{kbaufgabe}{Ich finde den Fehler beim Aufstellen eines Terms.}
\begin{kbblock}
\kbteil{a)}{Mia soll zu „Eine Zahl $x$ wird um 9 vermindert“ einen Term schreiben. Sie schreibt: $9 - x$. Finde den Fehler und schreibe den Term richtig.}{}
\item[]\kbraum{2}
\end{kbblock}
\begin{kbblock}
\kbteil{b)}{Jonas soll zu „Das Doppelte der Summe aus einer Zahl $x$ und 7“ einen Term schreiben. Er schreibt: $2 \cdot (x + 7)$. Prüfe, ob Jonas richtig übersetzt hat.}{}
\item[]\kbraum{2}
\end{kbblock}
\begin{kbblock}
\kbteil{c)}{Vier Schüler haben einen Text in einen Term übersetzt. Genau eine der vier Zeilen ist falsch. Welche? Schreibe den Term richtig.}{}
\kbfrage{(1) das Vierfache einer Zahl $x$, um 5 vergrößert: $4x + 5$ \\ (2) das Doppelte der Summe aus $x$ und 6: $2x + 6$ \\ (3) eine Zahl $x$, vermindert um 8: $x - 8$ \\ (4) die Hälfte einer Zahl $x$: $x : 2$}
\item[]\kbraum{2}
\end{kbblock}
\end{kbaufgabe}


% pflicht: terme-e1-k4-s2-v1, terme-e1-k4-s2-v2, terme-e1-k4-s2-v3
\begin{kbaufgabe}{Ich kann begründen, ob ein Term zum Text passt.}
\begin{kbblock}
\kbteil{a)}{Begründe, warum man für „das Doppelte der Summe aus $x$ und 5“ eine Klammer braucht. Setze dazu $x = 3$ in $2 \cdot (x + 5)$ und in $2x + 5$ ein.}{}
\item[]\kbraum{3}
\end{kbblock}
\begin{kbblock}
\kbteil{b)}{Entscheide bei jeder Aussage, ob sie wahr oder falsch ist. Begründe, ohne genau zu rechnen.}{}
\kbfrage{(1) „Eine Zahl $x$, vermindert um 3“ und „3 vermindert um eine Zahl $x$“ ergeben immer denselben Term. \\ (2) Das Doppelte einer Summe braucht immer eine Klammer. \\ (3) Es gibt eine Zahl $x$, für die $2 \cdot (x + 4)$ und $2x + 4$ denselben Wert haben.}
\item[]\kbraum{3}
\end{kbblock}
\begin{kbblock}
\kbteil{c)}{Emre sagt: „Für „das Doppelte der Summe aus $x$ und 6“ schreibe ich $2x + 6$. Verdoppeln heißt einfach: eine 2 davor.“ Begründe, ob Emre recht hat.}{}
\item[]\kbraum{3}
\end{kbblock}
\end{kbaufgabe}


% pflicht: terme-e1-k4-s3-v1, terme-e1-k4-s3-v2, terme-e1-k4-s3-v3
\begin{kbaufgabe}{Ich kann Terme und Termbäume einander zuordnen.}
\begin{kbblock}
\kbteil{a)}{Welcher Term gehört zum Termbaum?}{}
\kbdarunter{\termbaum{\cdot}{4}{\tb{+}{x}{6}}}
\kbantwort{\kbfeld}
\end{kbblock}
\begin{kbblock}
\kbanweisung{Fülle den Termbaum zum Term … aus.}
\kbteil{b)}{$2x - 7$}{}
\kbdarunter{\termbaum{}{\tb{}{}{}}{}}
\end{kbblock}
\begin{kbblock}
\kbteil{c)}{$(x - 2) \cdot 5$}{}
\kbdarunter{\termbaum{}{\tb{}{}{}}{}}
\end{kbblock}
\end{kbaufgabe}


% pflicht: terme-e1-k4-s4-v1, terme-e1-k4-s4-v2, terme-e1-k4-s4-v3
\begin{kbaufgabe}{Ich kann Terme für Aufgaben aus dem Alltag aufstellen.}
\begin{kbblock}
\kbteil{a)}{Ein Stromtarif kostet 11 € Grundpreis im Monat und 0,30 € je Kilowattstunde. Stelle einen Term für die Monatskosten in Euro bei $n$ Kilowattstunden auf. Was kostet ein Monat mit 150 Kilowattstunden?}{}
\item[]\kbraum{3}
\end{kbblock}
\begin{kbblock}
\kbteil{b)}{Ein Taxi kostet 4,50 € Grundpreis und 2,50 € je Kilometer. Stelle einen Term für den Fahrpreis in Euro bei $d$ Kilometern auf. Reichen 25 € für eine Fahrt von 9 km?}{}
\item[]\kbraum{3}
\end{kbblock}
\begin{kbblock}
\kbteil{c)}{Ein Handwerker berechnet 30 € für die Anfahrt und 45 € je Arbeitsstunde. Stelle einen Term für die Rechnung in Euro bei $t$ Stunden auf. Was kostet ein Einsatz von 3 Stunden?}{}
\item[]\kbraum{3}
\end{kbblock}
\end{kbaufgabe}
